\section{Countably many compatible partitions}\label{sec:countable}

The separation widths in \Cref{rounds:wide} may depend on the label.
We now use this freedom to construct the partitions. The analytic
construction is performed at finite stages. Its essential finiteness
property concerns each \emph{whole packet}, including all its infinite
product ends; local finiteness merely on compact subsets would not
suffice. The finite-label procedure is stated in \Cref{finite:cut}, and its detailed
proof is in \Cref{sec:packets,sec:cut,sec:local,sec:shield-global,sec:regularization}.
We adapt that proof using individual positive floors and finite
exceptional data on each whole packet.
This section establishes their hypotheses for the countable construction
and proves that those hypotheses persist.

\subsection{Initialization and the inductive data}

Let $D$, the centers $x_i$, and the pairs $(A_i,B_i)$ be as in
\Cref{rounds:wide}. We assume every $B_i$ is nonempty, since otherwise
$\diam_g M\le4D$. There is one label for each center, and an additional
label $0$ used first for radial confinement. The set of labels is finite
or countable; write it as $I=\{0,1,2,\ldots\}$, stopping the enumeration
if necessary. Choose
\begin{equation}\label{countable:floors}
 p_i^*>0,\qquad \sum_{i\in I}p_i^*=\frac1{20}.
\end{equation}
Use \Cref{rounds:wide} with
\begin{equation}\label{countable:widths}
 L_i=10+2\pi\sqrt{2/p_i^*}\qquad(i\ne0).
\end{equation}
We henceforth use the resulting complete metric $h\ge g$, its locally
finite list of log warps $p_\alpha$, and its stabilized scalar lower
bound $2/5$.

The initial weighted base measure is globally well defined:
\[
 \omega\,\dd\vol_h,\qquad
 \omega=\exp\Bigl(\sum_\alpha p_\alpha\Bigr)>0.
\]
The sum is locally finite. On a buffered compact set, only finitely many
nonunit factors $e^{p_\alpha}$ occur, so ordinary finite warped-product
formulas compute the stabilized scalar, weighted variations, and
invariant test forms there. These formulas agree on overlaps: adjoining
a unit factor changes none of them. In particular, the initial problem
is a weighted problem on the $n$-dimensional base $M$.

Fix numbers
\[
 0=t_0<t_1<t_2<\cdots<1,\qquad t_r\longrightarrow1,
\]
and use as many stages as there are labels. These are budget parameters,
unrelated to a spatial coordinate. Put
\begin{equation}\label{countable:constants}
 \begin{gathered}
 a_r=.60+.03t_r,\qquad b_r=1.12-.05t_r,
 \qquad \lambda_r=\frac14-\frac{t_r}{8},\\
 \Delta_r=t_{r+1}-t_r,\qquad
 e_r=\frac{\Delta_r}{800},\qquad
 \eta_r=\frac{\Delta_r}{200},\qquad k=399,
 \qquad C=\frac{4k}{k+1}=3.99.
 \end{gathered}
\end{equation}
Thus $\lambda_r\ge1/8$, $.60\le a_r<.63$, and
$1.07<b_r\le1.12$. No positive infimum of the $\Delta_r$ is required. In a stagewise
appendix lemma, its symbol $e_0$ denotes the current $e_r$.

We use the packet notation and regular splitting operation introduced in
\Cref{sec:finite-cuts}. The countable invariant below replaces the common
finite-label floor by the individual positive floors $p_j^*$.
At stage $r$, the processed labels are $0,\ldots,r-1$ and the remaining
set is $J_r$. A packet with $s$ hit entries has base dimension $m=n-s$.
Its continuous collapse map to $M$ is denoted by $q_E$. For each
$j\in J_r$ retain an extended-real continuous profile $\mu_j$ with
locally Lipschitz finite truncations, a smooth potential
$\Pi_j\ge p_j^*$, and disjoint open cores $U_j^+,U_j^-$ with
\begin{equation}\label{countable:profile-invariant}
 \overline{U_j^\pm}\subset\operatorname{int}\{\mu_j=\pm\infty\},
 \qquad 2|\grad\mu_j|\le a_r\mu_j^2+\Pi_j
 \quad\hbox{almost everywhere on the finite band}.
\end{equation}
For unprocessed center labels require
$q_E^{-1}(A_j)\subset U_j^+$ and $q_E^{-1}(B_j)\subset U_j^-$.
All data are compatible products in the packet ends. With
\begin{equation}\label{countable:required}
 G_j=E\setminus(U_j^+\cup U_j^-),\qquad
 D_E=\{x:\#\{j\in J_r:x\in G_j\}\ge m-1\},
\end{equation}
the scalar requirement is
\begin{equation}\label{countable:scalar-invariant}
 R_E\ge\lambda_r+b_r\sum_{j\in J_r}\Pi_j\quad\hbox{on }D_E.
\end{equation}
Here $R_E$ denotes the stabilized scalar expression, and integrals and
Laplacians use the product of the old warp lengths as weight. The
finiteness statement below makes these meanings unambiguous. A nonpositive
counting threshold means the whole packet.

\begin{lemma}[Initial profiles and radial blocks]\label{countable:initial}
The data in \eqref{countable:profile-invariant}--\eqref{countable:scalar-invariant}
exist at stage zero with $\Pi_i=p_i^*$. The radial forced signs can be
chosen so that every regular radial split preserving them has locally
finite compact wall and closed-side components before cylinder
attachment. On every compact radial block and buffer, only finitely
many center labels fail to be identically the negative pole with full
negative core, and only finitely many initial warps fail to be one.
\end{lemma}

\begin{proof}
For a positive number $p$, the decreasing crossing
\[
 2|y'|=.5y^2+p
\]
from $+\infty$ to $-\infty$ has length
\[
 T(p)=\int_{-\infty}^{\infty}\frac{2\,\dd y}{.5y^2+p}
      =2\pi\sqrt{2/p}.
\]
Extend it by its poles. For center $i$, begin this crossing at stretched
distance $2$ from $A_i$ and compose it with $\dist_h(A_i,\cdot)$.
The distance is $1$-Lipschitz, so the resulting profile obeys the
inequality with $.5$ in place of $a_0=.60$. Take cores
\[
 U_i^+=\{\dist_h(A_i,\cdot)<1\},\qquad
 U_i^-=\{\dist_h(A_i,\cdot)>T(p_i^*)+3\}.
\]
Their closures lie in the respective pole interiors. The stipulated
width puts $B_i$ strictly inside the negative core. In particular, on
the complement of the original closed ball $\overline B_g(x_i,2D)$,
the profile is the negative pole and the negative core is full.

There exists a smooth proper function $\rho:M\to\R$, bounded below,
with $|\rho-\dist_h(o,\cdot)|<1$ and $|\grad\rho|_h\le2$.
Here no uniform geometric assumption is necessary. To construct it,
use relatively compact coordinate balls in a locally finite cover,
convolve the distance on slightly larger balls, and blend with a
smooth partition of unity. The metric distortion and local convolution
errors can be made arbitrarily small on each ball. In differentiating
the blend, write each local approximation as the distance plus its
error; the terms involving derivatives of the partition sum to zero
on the common distance. Choose the value errors so that the remaining
partition-derivative terms have sum less than a fixed small number at
each point. A locally finite refinement and successively smaller
positive errors achieve this, while keeping the total value error
below one. Properness follows from the bounded value error and
properness of the complete metric $h$.

Choose alternating positive and negative plateaus in the $\rho$-axis,
with transition intervals wider than $2T(p_0^*)$. Rescale the preceding
crossing to these intervals, reversing its sign on alternate intervals.
The gradient bound for $\rho$ gives the profile inequality with
coefficient $.5$ and potential $p_0^*$. Choose the core closures
strictly inside the plateaus. Regular values in the remaining margins
give compact smooth barrier annuli, and a compact central barrier,
whose closures stay in the same pole plateaus. Their value bands tend
to infinity and are locally finite.

A wall respecting the signs cannot cross a separating plateau band;
each component is closed in one compact radius block and hence compact.
A component of either closed side is trapped between two separator
bands of the opposite sign, or in the initial compact block.
Regularity of the defining function gives manifold-with-boundary
charts for these sides. A compact block therefore contains only
finitely many components. The same charts prove local finiteness of
the component families. When $M$ is compact, use a radial label with
constant positive pole and full positive core; its split has no wall
and its only side is compact.

The scalar requirement initially holds everywhere because
\[
 \lambda_0+b_0\sum_i p_i^*=.25+1.12/20=.306<.4.
\]
Finally a compact block with a compact buffer meets only finitely many
original closed outer balls and finitely many initial warp supports.
Every other center label and warp has the asserted constant data there.
\end{proof}

\subsection{Finite exceptional data on entire packets}

A \emph{downward diagram} of a packet consists of it and all its further
faces. An \emph{upward diagram} consists of it and all its cofaces.
Each is finite at a fixed stage, by \Cref{packet:split}. A finite union
of either type is finite. The constructions below make finitely many
directed passes through such diagrams. They never ask for closure
under arbitrary alternating face and coface steps.

Radial confinement gives the geometric reason to expect finite data on
a whole packet. At the first split, a compact radial side or wall
component is prolonged by cylinders; its entire collapse image stays
in that original compact component. At every later split a child $A$
of $E$ has
\begin{equation}\label{countable:collapse-recursion}
 q_A=q_E\circ p_A,\qquad
 p_A=\begin{cases}
 \text{inclusion},&\text{on the wall or bulk},\\
 \text{projection to the wall},&\text{on a fresh cylinder}.
 \end{cases}
\end{equation}
Hence $q_A(A)\subset q_E(E)$, including in all simultaneous normal
directions. A finite incidence diagram therefore has image in a finite
union of compact original radial blocks. Only finitely many original
outer balls and initial nonunit warp supports meet that image.
Shielding changes no collapse map, and regularization chooses new
walls inside old parents, so both preserve this confinement.

This argument controls the initial supports. Later smooth data need
not factor through $q_A$: they also use data copied from faces. We
must therefore retain finite exceptional lists through those directed
dependencies, as well as preserve every other entry exactly.

\begin{definition}[Finite exceptional data]\label{countable:finite-data}
At stage $r\ge1$, on each whole packet $E$ there are finite sets
$I_E\subset J_r$ of label indices and $F_E$ of old warp indices such
that, for $j\notin I_E$,
\begin{equation}\label{countable:passive}
 \mu_j\equiv-\infty,\qquad U_j^-=E,\qquad U_j^+=\varnothing,
 \qquad \Pi_j\equiv c_rp_j^*.
\end{equation}
For every old warp index outside $F_E$, its length is identically one.
The number $c_r$ is the same on all packets and is prescribed by
\begin{equation}\label{countable:passive-multiplier}
 c_0=1,\qquad c_{r+1}=(1+3\eta_r/2)c_r.
\end{equation}
For any finite diagram used in a comparison, take the union of its
finite exceptional sets, adding the current label explicitly. This
one finite union then works simultaneously on all its whole packets.
At stage zero the analogous property is required on every compact
radial block and buffer, with $c_0=1$.
\end{definition}

The sets $I_E$ need not consist exactly of the nonconstant entries: an
entry may be retained in the exceptional list even if it happens to be
constant. This convention avoids any need to discard previous exceptions.
Only the existence of a finite list matters.

\begin{lemma}[Consequences and directed dependencies]\label{countable:locality}
Suppose the packet geometry and \Cref{countable:finite-data} hold at a
finite stage. Then the following assertions hold.
\begin{enumerate}[label=(\roman*),leftmargin=*]
\item The sum of potentials on a whole packet is a smooth function
with bounded derivatives of every fixed order. The undecided counts,
required sets, and guards use only a finite list and are closed.
\item On each whole packet and each finite comparison diagram, the old
weight, stabilized scalar, stability forms, and weighted Laplacian are
computed with one finite list of warps, after padding with unit factors.
The bounded geometry and polynomial weighted growth conclusions of
\Cref{packet:geometry} apply, including on any finite Euclidean product.
\item Compatible choices with positive packetwise tolerances can be
made in the following order: fix individual requests, propagate numerical
requests down the finite coface chains, and then construct data upward
from deeper faces. Enlarging a coface's exceptional list by the union
over its downward source diagram suffices for the last pass. No data
are propagated back from a newly constructed coface to an already
constructed face.
\end{enumerate}
\end{lemma}

\begin{proof}
On $E$, the sum of potentials is exactly
\begin{equation}\label{countable:potential-sum}
 \sum_{j\in J_r}\Pi_j
 =\sum_{j\in I_E}\Pi_j+c_r\sum_{j\in J_r\setminus I_E}p_j^*.
\end{equation}
The second term is a finite spatial constant. Indeed
\[
 1\le c_r\le\exp\Bigl(\frac32\sum_{s<r}\eta_s\Bigr)
 \le e^{3/400}.
\]
The first term is a finite sum; its derivative bounds follow from
compatible smoothness and the compact-part/product-end decomposition
in \Cref{packet:geometry}. On a finite comparison diagram choose common
product thresholds for its finite exceptional list. Every other entry is already a spatial
constant with the same value on each copy. The full sum is therefore
an exact face pullback beyond those thresholds; no supremum of
infinitely many thresholds is taken. Passive entries have empty
undecided set. Thus a set requiring at least $a\ge1$ undecided labels
is the finite union, over $a$-element subsets of $I_E$, of the intersections of
their closed $G_j$'s. This proves closedness, including for the guards
with threshold $m-2$. Threshold zero gives $E$.

For (ii), take the union of the finitely many relevant old warp
indices on the finite diagram and set every missing entry to one.
Writing $u_\alpha=\log f_\alpha$ and $U=\sum_\alpha u_\alpha$, the
scalar expression is
\[
 R_E=\Scal_{h_E}-2\Delta_{h_E}U-|\dd U|_{h_E}^2
                        -\sum_\alpha|\dd u_\alpha|_{h_E}^2.
\]
A padded logarithm is zero; the product weight is unchanged. The
weighted Laplacian is
$\Delta_{h_E}+\langle\dd U,\dd\,\cdot\,\rangle_{h_E}$ and is
unchanged as well. Weighted perimeter and its first and second
variations depend on the same weight. Equivalently, a flat unit
factor adds zero second fundamental form and zero curvature to the
invariant variation calculation. Thus the wall and seam forms also
agree. The finite-diagram geometric lemma now gives completeness,
bounded geometry, positive lower and finite upper bounds for the
weight, and center-uniform polynomial volume growth. It gives the
same statements after taking a finite Euclidean product. The
regularity problem for a minimizing phase remains on the base, with
singular dimension at most $m-8$ by \Cref{cut:gmt}; padding does not
alter that dimension.

For (iii), a particular face receives requests from only its finitely
many cofaces and coface chains. Allowing for the finitely many fixed
product and cutoff bounds, the minimum of the requested positive
errors is positive. Settle all these numbers before constructing the
approximations. A request only decreases an error tolerance: it does
not copy any coface value into the face. During the subsequent upward
pass, the inputs on a packet are its own data and already constructed
copies from its finite downward diagram. Their exceptional sets have
a finite union. Outside that union every copy in
\eqref{countable:passive} agrees with every other copy. A newly
constructed coface is not used to reconstruct its sources. The same
argument applies to finitely many successive directed passes, and
proves the stated dependency rule.
With these numerical tolerances fixed, at each upward pass the output
on $E$ depends only on its original data and the outputs on its strict
faces. Its source set for copied functions is therefore contained in
the finite downward diagram of $E$. The preceding downward pass changes only positive
numbers, not the input functions. In particular two packets sharing
a face do not copy one another's newly constructed data through that face.
\end{proof}

\subsection{One countable invariant step}

\begin{proposition}[Preservation of the invariant]\label{countable:step}
Suppose that all nonempty parents at stage $r$ have base dimension at
least two and satisfy \eqref{countable:profile-invariant}--\eqref{countable:scalar-invariant}.
At $r=0$ use the compact radial blocks of \Cref{countable:initial}; at
$r\ge1$ suppose also that the packet geometry and
\Cref{countable:finite-data} hold. Processing label $r$ produces the
compatible smooth split, preserves the forced signs, and gives the
same assertions at stage $r+1$. Exactly $399$ new warp factors, with
common indices, are appended at this stage. Each whole resulting
packet has a finite relevant old warp list.
\end{proposition}

\begin{proof}
We verify the choices in their order of dependence.

\emph{The current pressure and minimizing phases.}
The pressure construction of \Cref{cut:pressure} uses a positive floor
only for the label being processed. Replace its finite-family floor
by $p_r^*$. Compatible clipping heights can first be chosen above the
local barrier bounds, with a uniform extra barrier margin of one.
Then use local smoothing errors to obtain
\[
 2|\grad p_0|\le .65p_0^2+1.01\Pi_r.
\]
All smoothing requests involve only this one profile and finite
incidence diagrams. The elementary inequality
\[
 .67(y+c)^2-.65y^2\ge-21.775c^2
\]
shows that one common nonempty shift interval, for example
\begin{equation}\label{countable:shift}
 |c|<\min\{.01\sqrt{p_r^*},1/2\},
\end{equation}
preserves both $2|\grad(p_0+c)|\le .67(p_0+c)^2+1.02\Pi_r$ and
the barrier inequalities. Floors for other labels do not enter this
interval.

After the radial stage, \Cref{countable:locality} supplies all bounded
geometry and weighted polynomial-growth hypotheses of
\Cref{cut:generic-uniqueness}. There are countably many packets,
Euclidean product dimensions, and tests in that lemma's exhaustion.
Remove the countable exceptional shift sets and choose a single $c$
in \eqref{countable:shift}. The resulting true minimizers have the
common product limits of \Cref{cut:end-convergence}; choose them before
any shielding update. At the initial radial step no uniqueness or
polynomial growth at spatial infinity is used. The smooth global
positive weight and locally bounded smooth pressure give local BV
existence and strict compactness by \Cref{cut:bv-compactness}. The
relatively compact radial barriers force the phases by compact
comparisons. All variations and tests lie in buffered compact windows
with finite warp lists, so the formulas in \Cref{cut:raw-domination}
apply. The GMT and shielding dimensions are the base dimensions at
this first step as well.

\emph{Anchors, guards, and shielding.}
On each post-radial packet retain its finite exceptional list, enlarged
by its finite downward source diagram and by the current label.
Labels outside it are already forced everywhere and require no shell
addition. Choose compatible smooth clipped anchors as in
\Cref{local:anchors}, retaining the exact value $-1$ for a passive
profile. For each packet of dimension $m\ge8$ and each set of $m-7$
remaining labels, the image of its singular set under the corresponding
smooth anchor map has $(m-7)$-dimensional Lebesgue measure zero.
There are countably many such conditions, even if all finite subsets
of the countable label set are included. In the countable product of
uniform probability intervals $(-.3,.3)$, each bad cylinder has
measure zero by its finite-dimensional marginal. A threshold vector
outside their union therefore exists. At most $m-8$ anchors equal
their thresholds at any singular point. The sign exclusions at poles
follow from the $.1$ anchor approximation, exactly as in
\Cref{local:anchors}.

Process the finitely many face depths from deepest upward, keeping the
true cuts fixed. Suppose all strict faces of $E$ have been shielded.
Its remaining risky set
\[
 K_E=\Sigma_E\cap
 \{x:\#\{j\in J_r\setminus\{r\}:x\in G'_j\}\ge m-2\}
\]
is closed. It is compact: an escaping sequence would, by
\Cref{packet:escape}, have a product limit over a further face with
compact convergent base coordinates. The finite exceptional list and
closed core complements pass the undecided count to that limit.
The face has a smaller dimension, so the limit lies in its guard.
The product convergence of true cuts forces the limiting point to be
singular, since regular-point graphical convergence excludes nearby
singular points. This contradicts the completed shielding of that
face. This argument also treats escape in several normal coordinates.

Apply \Cref{local:flat-pieces,local:shell,local:sharp} on this compact
risky set, only for the finite list which can still be undecided.
Enumerate the signed compact flat pieces. Choose each shell scale by
\Cref{shield:maximum}, hold it fixed, and then send its coface fronts
far enough out by \Cref{shield:extension}. In the cap construction
of \Cref{shield:caps}, the crossing for label $j$ satisfies
\[
 2|w'|\le .58w^2+p_j^*.
\]
Its length is finite; the inequality is paid by $\Pi_j\ge p_j^*$.
The minimum-of-branches formula in that lemma adds no simultaneous
normal derivative cost. Its cutoff finishes before the cap crossing
starts, so whenever the shell branch is active the whole shell
potential is available. Opposite signs remain in disjoint strict
anchor regions, and cores only enlarge.

Only finitely many additions are required on $E$. Indeed, every point
of the initial compact $K_E$ eventually belongs to an open forced core
for each of its strict anchors among the finite exceptional list.
At most $m-8<m-2$ labels can then remain undecided there. The residual
risky subsets of the initial $K_E$ are nested and closed, and their
intersection is empty. Compactness gives a finite index at which the
residual is empty. Stop then. Each specified coface receives additions
from finitely many source faces through finitely many chains. Thus
only finitely many additions and finite minima or maxima affect any
whole packet. In particular no infinite limiting profile is taken on
that packet.

To control all the countably many packets at once, enumerate each
possible local addition and assign positive allowances with total less
than $.12$. Split each allowance between its shell-scale choice and
its simultaneous finite coface extension choices. Starting the
bookkeeping form bound at $.78$, the maximum estimate and the
extension estimate give one $\theta_*<.9$ such that
\begin{equation}\label{countable:shield-form}
 \int1.16\sum_{j\ne r}D_jv^2
 \le\theta_*\bigl(\mathcal N(v)+e_r\|v\|_2^2\bigr).
\end{equation}
Here $\mathcal N$ has coefficient $1.7$ on a wall and coefficient
$.8$ on bulk plus $1.7$ on its attached cylinder. These are the
same coefficients on every packet and product. Only a finite list of
additions occurs on each post-radial packet, so the bound persists
after a whole depth has been treated.

Initially, instead work in separated compact transition shells with
compact buffers between radial barriers. Each has finitely many
exceptional labels and warps. Its risky set is compact, and the same
finite termination applies; there is no coface extension yet. Choose
supports within the shell buffers. Their local finiteness makes all
updates locally finite on $M$, and a compact test meets only finitely
many of them. The summable allowances prove
\eqref{countable:shield-form} for every such test. Every compact radial
side block meets only finitely many shell buffers. This supplies the
whole-packet finite lists after the first attachment.

\emph{Counts and the quantitative transfer.}
Core enlargement gives $G'_j\subset G_j$. Only the finite exceptional
list can contribute to the counts on a post-radial packet; the analogous
statement holds on buffered compact blocks at the radial step.
Thus \Cref{packet:counts} places every new required set in both the old
required set and the shielded guard under projection, including along
attached cylinders. The finite lists and closed core complements also
pass these count inclusions to all further-face product limits.

The guarded tube and comparison constructions of
\Cref{regular:guards,regular:partition} use these closed guards,
regular graphical convergence on product limits, and finite
packetwise derivative bounds. All are now available. Numerical
errors are allocated down coface chains before the defining functions
are constructed upward.

Apply the quantitative transfer of \Cref{regular:transfer} with the
current $t_r$, $\Delta_r=t_{r+1}-t_r>0$, and $e_0=e_r=\Delta_r/800$.
Its stage margins \eqref{regular:eta} were proved for these general
adjacent parameters, so they do not require equally spaced stages.
The potential calculation also applies: on each finite comparison
diagram, \eqref{countable:potential-sum} expresses the old potential
sum as finitely many smooth functions plus one spatial constant.
Only finitely many shielding additions occur there. Thus, with
$\widetilde\Pi_j=\Pi_j+D_j$, the same raw form and comparison estimates
control
\begin{equation}\label{countable:raw-target}
 \lambda_{r+1}+(1+3\eta_r)b_{r+1}
                         \sum_{j\ne r}\widetilde\Pi_j
\end{equation}
near the new required set, with the positive mass and gradient reserves
proved there. The constant tail has no derivative error; potential
approximation errors concern only the finite exceptional list. At the
radial step these estimates are made on buffered compact blocks, as
in that proof.
The seam and potential smoothing can consequently retain an $e_r$
local mass gap. This positive number is fixed for the current stage;
all geometric tolerances may depend on the packet. No infimum over
future stages or over all packets is taken.

\emph{Exact passive constants in regularization.}
Restrictions, bulk inclusions, and cylinder projections preserve
identically constant functions. Seam smoothing of an old identically
unit warp leaves it one. For the remaining potentials, smooth
$(1+3\eta_r/2)\widetilde\Pi_j$ as in
\Cref{regular:profiles,regular:constant-preservation}, so that
\begin{equation}\label{countable:smooth-potentials}
 (1+\eta_r)\widetilde\Pi_j\le\Pi_j^{\rm new}
                     \le(1+2\eta_r)\widetilde\Pi_j.
\end{equation}
On a finite exceptional list, allocate the required relative errors
separately, using $p_j^*>0$ for its $j$th entry. If an aggregate error
is prescribed, subdivide it among that finite list. If an entry is
passive on the whole source diagram, set its output \emph{exactly} to
$c_{r+1}p_j^*$. This is compatible with the smoothing procedure:
local convolution of a constant is that constant, partition weights
sum to one, and every copied end value is the identical constant.
Equivalently this entry can simply be omitted from all smoothing and
restored as that constant afterward. We only impose this rule on
whole constant inputs and their agreeing copies, so it makes no claim
about preserving the edge of an arbitrary nonsmooth plateau.

Before blending on a coface, enlarge its exceptional list by the union
over its entire finite downward source diagram. Thus a face-active
entry is recorded even if the coface's individual input was passive.
The remaining copies agree identically. Coface requests merely adjust
numeric tolerances; they do not introduce new source values on faces.
This proves exact passive reproduction on the completed diagram and
precludes a cascade through infinitely many same-depth neighbors.
At the radial step choose a compact buffer for each compact radial
block, enlarged to contain the finitely many chart closures and shell
supports used by its approximations. Only finitely many original
outer balls and initial nonunit warp supports meet this buffer.
Enlarge the exceptional list by precisely these indices. Every other
entry is the same constant on every averaging neighborhood and is
therefore reproduced exactly before cylinder attachment.

The local-slope argument of \Cref{regular:profiles} justifies restriction
of profiles to the wall: an almost-everywhere gradient bound by a
continuous function gives the corresponding local Lipschitz slope
bound at every point of each finite band. Restriction and cylinder
projection then preserve it. The small metric change increases slopes
by at most $1+\eta_r$, and the first inequality in \eqref{regular:eta}
together with \eqref{countable:smooth-potentials} gives the next profile
inequality.
Continuous pullback and core enlargement preserve all pole-neighborhood
and forced-preimage conditions.

\emph{Corrections, supersolutions, and the new warp group.}
The sum in \eqref{countable:potential-sum} is now smooth, bounded, and
exactly compatible, as are its finite-order derivatives. On every
required finite diagram, pad the old warps by units. In a product end
over $B$ the operator of \Cref{regular:correction} therefore satisfies
exactly
\[
 L_A=L_B-C\Delta_{\rm normal},\qquad
 L_A=-C\Delta_{\rm old}+R_{\rm old}-V_A^*+K_A,
 \qquad
 V_A^*=\lambda_{r+1}+b_{r+1}\sum_{j\ne r}\Pi_j^{\rm new}.
\]
In particular the potential formula does not depend on the base
dimension or on how many unit factors were padded. The correction is
zero near the new required set and produces a positive form gap on
each whole packet. The proof of \Cref{regular:correction} uses finite
minima of the gaps on the packet's face diagram. Its bounded resolvent
and the compatible positive supersolutions in
\Cref{regular:resolvent,regular:solutions} apply with the packetwise
bounded geometry and polynomial weighted volume already proved.
Their tolerances are chosen after the operators and their individual
exact solutions are fixed, and before the final upward blending pass.
Thus their finite-dependency hypothesis is precisely (iii) of
\Cref{countable:locality}.

Let $v_A>0$ be the resulting supersolution with $L_Av_A\ge1/2$.
Append the common group of $k=399$ circle lengths
$f=v_A^{2/(k+1)}=v_A^{1/200}$. The scalar identity in
\Cref{regular:step} gives
\[
 R_{\rm new}=R_{\rm old}-C\frac{\Delta_{\rm old}v_A}{v_A}
           =V_A^*-K_A+\frac{L_Av_A}{v_A}
           \ge V_A^*\quad\hbox{on }D_A.
\]
Only the number of \emph{new} factors enters the coefficient $C$.
The old finite list may vary from packet to packet; unit padding has
no effect. Each stage adds just this one finite group, with the same
indices in every status. All its lengths are positive and bounded
above and away from zero on a whole packet, and exactly compatible.
The initial lists and their finite propagated exceptions remain finite,
so the complete old list at stage $r+1$ is still finite on each whole
packet and each finite comparison diagram.

Finally the regular split preserves compact truncations and finite
incidence by \Cref{packet:split}. Initially the radial separators make
all components compact before prolongation. Compatible smooth metrics
and finite relevant warps then give the complete packet geometry by
\Cref{packet:geometry}. The recursion
\eqref{countable:collapse-recursion} preserves confinement in original
radial blocks. All exceptional-list enlargements just described are
finite, and all passive constants have multiplier $c_{r+1}$.
This completes the induction step, including the first radial step.
\end{proof}

\subsection{Termination in dimension and finite-stage walls}

\begin{lemma}[The one-dimensional obstruction]\label{countable:dimension-one}
No nonempty one-dimensional packet can satisfy the invariant at a
finite stage of the construction.
\end{lemma}

\begin{proof}
Such a packet is complete by \Cref{packet:geometry} and has only
finitely many relevant positive smooth warps $f_1,\ldots,f_q$ on
its whole base. Its arclength universal cover is $\R$, including
when the packet itself is a circle. On this line put
$u_\alpha=\log f_\alpha$ and $P=\sum_{\alpha=1}^q u_\alpha$.
The warped-product scalar formula reduces to
\[
 R=-2P''-(P')^2-\sum_{\alpha=1}^q(u_\alpha')^2.
\]
The required count in dimension one is zero, so the invariant holds
everywhere and gives $R\ge\lambda_r>0$. Consequently
\[
 -2P''\ge\lambda_r+(P')^2,\qquad
 \frac{\dd}{\dd s}\arctan\frac{P'}{\sqrt{\lambda_r}}
 \le-\frac{\sqrt{\lambda_r}}2.
\]
Integration over a line interval longer than
$2\pi/\sqrt{\lambda_r}$ contradicts the range of the arctangent.
If $q=0$, the scalar formula is simply $R=0$. The argument takes
place at one finite stage and uses an ordinary finite sum.
\end{proof}

\begin{proposition}[The countable compatible construction]\label{countable:construction}
All labels can be processed successively. At every finite stage the
packet diagram has the invariant above and no nonempty packet of
dimension at most one. After the radial stage it has finite exceptional
data on whole packets; at stage zero this assertion holds on buffered
compact radial blocks.
Its finite reconstruction $X_r$ has a proper collapse $q_r:X_r\to M$
and closed signed sides $P_i^{(r)},N_i^{(r)}$ for each processed center
label, with
\[
 P_i^{(r)}\cup N_i^{(r)}=X_r,\qquad
 P_i^{(r)}\cap N_i^{(r)}=W_i^{(r)},
\]
where $W_i^{(r)}$ is its wall. The forced preimages satisfy
\[
 q_r^{-1}(A_i)\subset P_i^{(r)}\setminus W_i^{(r)},\qquad
 q_r^{-1}(B_i)\subset N_i^{(r)}\setminus W_i^{(r)}.
\]
These inclusions hold in open phases and persist at every subsequent
finite stage. No $n-1$ distinct walls, even including the radial wall,
have a common point. Each center wall lies over its original open
outer ball, and its components are compact and locally finite.

Old horizons can be enlarged by identifications preserving collapse,
all walls, and all signed statuses. For any prescribed sequence
$\eps_r>0$ with $\sum_r\eps_r<1$, the next finite collar compression
can be chosen with collapse displacement less than $\eps_r$, preserving
the previously decided wall strata and signed sides. These are
finite-stage conclusions; the limiting fixed-copy construction is
made in the next section.
\end{proposition}

\begin{proof}
Initialization is \Cref{countable:initial}. If the construction has
reached stage $r$, its parents have dimension at least two: a first
nonempty output of smaller dimension would have dimension exactly
one, since a hit decreases dimension by one and a sign preserves it.
\Cref{countable:step} produces the next invariant, and
\Cref{countable:dimension-one} excludes that first output. Thus the
induction continues through every finite stage. There is no passage
to an infinite-stage analytic metric, weight, or torus.

Apply \Cref{packet:reconstruction} to each finite initial segment.
Its hypotheses are the regular compatible splits and the compact
radial components already established. Its status pieces are exactly
the corresponding compact packet truncations. If $n-1$ walls met,
assign those labels hit status and choose a closed sign status for
each other processed label containing the point. The corresponding
nonempty truncation would have dimension one, a contradiction.
Forced signs persist under the continuous collapse recursion and
enlargement of the cores. In particular
\[
 q_r(W_i^{(r)})\subset M\setminus(A_i\cup B_i)
                         \subset B_g(x_i,2D).
\]
Properness of $q_r$ and compactness of the original closed outer ball
make each wall component compact; regular local wall charts give
local finiteness, as in \Cref{packet:reconstruction}.

The horizon identifications in \Cref{packet:reconstruction} preserve
collapse, walls, and signed statuses. Its collar compressions preserve
the old strata and their sides, and \Cref{packet:small-displacement}
allows their collapse oscillations to be bounded by the prescribed
$\eps_r$. These choices are made separately on the locally finite
compact component family; no uniform collar width or horizon at
infinity is required. The next section chooses these finite-stage
homeomorphisms coherently on one fixed copy of $M$.
\end{proof}
